\clearpage\fancyhead[L]{\small JUNO.}\hypertarget{native-335-lower}{}\label{native-335-lower}
\begin{center}\textbf{JUNO.}\par\vspace{4pt}\hrule\vspace{4pt}\fontsize{10}{12}\selectfont Monatliche Correspondenz. Band XXIV. Seite 186\,..\,188. 1811. August.\par\vspace{4pt}\hrule\vspace{4pt}\end{center}

\par\indent Im Septemberheft 1808 dieser Zeitschrift [S. 307 d. B.] haben wir unseren Lesern die Elemente der \emph{Juno} mitgeteilt, die Professor \textsc{Gauss} zum achten Mal verbessert hatte. Seither ließ sich eine weitere Verbesserung kaum vornehmen, da die fünfte Opposition der \emph{Juno}, die im Januar 1810 eintrat, nirgends beobachtet wurde. Umso wichtiger war es daher, die Opposition im April dieses Jahres nicht zu versäumen. Der Planet war ungewöhnlich lichtschwach, was die Beobachtung besonders schwierig machte. Wahrscheinlich ist Professor \textsc{Gauss} der Einzige, dem die Beobachtung dieser Opposition gelungen ist. Hier war der Planet selbst im so lichtstarken Meridianfernrohr der Seeberger Sternwarte nie sichtbar. Umso wertvoller sind die wenigen Beobachtungen, die Professor \textsc{Gauss} bei der äußerst geringen Helligkeit des Planeten nur mit dem Kreismikrometer anstellen konnte. Zum Vergleich dienten $\varphi$ Librae und einige andere nahe gelegene Sterne der \emph{Histoire céleste}. Wenn ihre Positionen künftig genauer bestimmt sind, werden die folgenden Ergebnisse noch eine kleine Korrektur brauchen.
\Needspace{10\baselineskip}\begin{center}\fontsize{10}{12}\selectfont\setlength{\tabcolsep}{4pt}\renewcommand{\arraystretch}{1.15}
\begin{tabular}{l|c|c|c}
1811 & \shortstack{Mittlere Zeit\\in Göttingen} & \shortstack{Rektaszension} & \shortstack{Südliche\\Deklination} \\
\hline
April 22 & $9^{\mathrm u}\,51^{\mathrm m}\,35^{\mathrm s}$ & $216^\circ\,41'\,50''1$ & $0^\circ\,58'\,16''$ \\
\phantom{April }24 & $10\;\;32\;\;55$ & $216\;\;17\;\;58.7$ & $0\;\;46\;\;50$ \\
\phantom{April }25 & $10\;\;19\;\;22$ & $216\;\;6\;\;3.6$ & $0\;\;40\;\;4$ \\
\phantom{April }26 & $10\;\;17\;\;30$ & $215\;\;54\;\;7.3$ & $0\;\;33\;\;51$ \\
\end{tabular}
\end{center}

\SourceLeaf{336}\Needspace{12\baselineskip}\par\indent Die Deklination vom 22. April ist zweifelhaft, und auch die vom 24. ist weniger zuverlässig als die beiden folgenden, bei denen der Planet günstiger stand. Aus diesen Beobachtungen hat Professor \textsc{Gauss} das folgende Ergebnis für die Opposition gewonnen.
\begin{center}\fontsize{10}{12}\selectfont\setlength{\tabcolsep}{4pt}\renewcommand{\arraystretch}{1.15}
\begin{tabular}{l}
$\text{\AstroSymbols ☍}$ 1811. 24. April $19^{\mathrm u}\,20^{\mathrm m}\,12^{\mathrm s}$ Mittlere Zeit in Göttingen \\
$214^\circ\,8'\,48''3$ wahre heliozentrische Länge \\
$12\;\;55\;\;2.0$ wahre geozentrische Breite, nördlich \\
\end{tabular}
\end{center}


\Needspace{20\baselineskip}\par\indent Die starke Abweichung dieser Position von den Elementen, die auf den ersten vier Oppositionen beruhen, zeigt nun auch bei Juno den Einfluss der Störungen, die besonders \emph{Jupiter} verursacht. Für Juno wird ihre Berechnung einen ebenso ungeheuren Aufwand erfordern wie für Pallas. Es ist nur angemessen, dass diese Arbeit zuerst für die früher entdeckte Pallas abgeschlossen wird. Deshalb hat sich Professor \textsc{Gauss} vorerst damit begnügt, neue elliptische Elemente auf die Oppositionen von 1806, 1807, 1808, 1811 zu stützen. Wir geben sie hier wieder.

\begin{center}\fontsize{10}{12}\selectfont\setlength{\tabcolsep}{4pt}\renewcommand{\arraystretch}{1.15}\begin{tabular}{llr}
\shortstack[l]{Epoche der mittleren Länge 1811\\Meridian von Göttingen} & .\;.\;.\;.\;. & $177^\circ\,48'\,1''8$\\
tägliche mittlere tropische Bewegung &  & $813''2486$\\
Länge des Perihels 1811 & .\;.\;. & $53^\circ\,14'\,32''4$\\
Länge des aufsteigenden Knotens 1811 &  & $171^\circ\,9'\,13''5$\\
Neigung der Bahn & .\;.\;.\;.\;.\;.\;. & $13\;4\;27.0$\\
Exzentrizität & .\;.\;.\;.\;.\;. & $=\mathrm{Sin}\;14\;44\;9.1$\\
Logarithmus der großen Halbachse &  & $=0.4265711.$\\
\end{tabular}\end{center}

\par\indent Nach diesen Elementen hat Herr \textsc{Wachter}, ein geschickter Schüler von Professor \textsc{Gauss}, die folgende Ephemeride für den Lauf der \emph{Juno} im Jahr 1812 berechnet. [vom 23. Februar bis 5. November]
